\begin{abstract}
We introduce \textsc{QLoRA}, a memory-efficient finetuning technique enabling the adaptation of large language models containing up to 65B parameters using a single 48GB GPU, while maintaining performance parity with full 16-bit finetuning. In \textsc{QLoRA}, gradients are propagated through a 4-bit quantized, frozen backbone model into trainable Low Rank Adapters~(LoRA). Our leading model suite, dubbed \textbf{Guanaco}, achieves superior results over all prior publicly available models on the Vicuna benchmark, attaining 99.3\% of ChatGPT's performance level after only 24 hours of finetuning on one GPU. Key contributions of \textsc{QLoRA} to memory optimization without loss of effectiveness include: (a) the introduction of 4-bit NormalFloat (NF4), a novel data format optimized for normally distributed parameter weights; (b) the application of Double Quantization, which compresses the quantization constants themselves to further reduce memory consumption; and (c) Paged Optimizers, which mitigate memory usage spikes during training. Utilizing \textsc{QLoRA}, we have finetuned over 1,000 models, conducting a thorough assessment of instruction-following and chatbot capabilities across 8 instruction datasets, different architectures (such as LLaMA and T5), and model sizes that are typically impractical for standard finetuning methods (including models with 33B and 65B parameters). Our findings indicate that QLoRA-based finetuning on small, high-quality datasets yields results that surpass previous state-of-the-art, even when utilizing models of reduced scale. We offer a comprehensive evaluation of chatbot performance through both human raters and GPT-4-based assessments, demonstrating that GPT-4 judgments provide a cost-effective and adequate proxy for human evaluations. Moreover, our study suggests that prevalent chatbot benchmarks do not reliably reflect chatbot capabilities. A targeted lemon-picked analysis is presented to illustrate failure cases of \textbf{Guanaco} when compared to ChatGPT. All models and accompanying code, including specialized CUDA kernels for 4-bit training, are released for public use.\footnote{\url{https://github.com/artidoro/qlora} and \url{https://github.com/TimDettmers/bitsandbytes}}
\end{abstract}

\section{Introduction}

Adapting large language models (LLMs) through finetuning has been established as a robust approach for enhancing their capabilities \citep{min2021metaicl, wei2021finetuned, ouyang2022training, wang2022super, wang2022self, liu2022few} and for incorporating or removing specific model behaviors \citep{ouyang2022training,askell2021general,bai2022training}. Nevertheless, the computational demands for finetuning especially large models are considerable; for example, standard 16-bit finetuning of a LLaMA model with 65 billion parameters~\citep{touvron2023llama} consumes over 780 GB of GPU memory. Although advances in quantization can compress LLM memory usage for inference \citep{dettmers2022llm, dettmers2022case, frantar2022gptq, xiao2022smoothquant}, these benefits are largely restricted to inference phases, as such strategies fail to extend to training scenarios \citep{wortsman2023stable}.

Here, we show for the first time that finetuning a 4-bit quantized model can be accomplished without any loss in model effectiveness. Our proposed approach, \textsc{QLoRA}, introduces an innovative, high-precision quantization strategy that reduces pretrained model parameters to 4 bits, then supplements them with a compact set of tunable Low-rank Adapter weights \citep{hu2021lora}, 

which are updated by backpropagating gradients through the quantized parameter space.

Employing \textsc{QLoRA} cuts the average GPU memory usage required to finetune a 65B parameter model from over 780GB down to less than 48GB, all while preserving both training efficiency and predictive quality relative to a fully finetuned 16-bit reference. This represents a major leap in making LLM finetuning accessible: now, state-of-the-art models of unprecedented scale are within reach for finetuning on just a single GPU. Leveraging \textsc{QLoRA}, we have created the \textbf{Guanaco} model suite, with the second largest configuration achieving 97.8\% of ChatGPT’s performance on the Vicuna~\citep{vicuna2023} benchmark, trainable in under 12 hours using a single consumer GPU; scaling to the largest Guanaco model, we reach 99.3\% with one professional GPU in 24 hours, effectively closing the performance gap with ChatGPT on this benchmark. Notably, our smallest \textbf{Guanaco} model (7B parameters) needs only 5 GB of memory and surpasses the 26 GB Alpaca model by over 20 percentage points in Vicuna benchmark accuracy (see Table~\ref{tbl:vicuna_eval}).

\textsc{QLoRA} presents several advancements aimed at minimizing memory requirements while maintaining high performance: (1) \textbf{4-bit NormalFloat}, a quantization format specifically tailored for normally distributed variables that is theoretically optimal in terms of information retention, and empirically surpasses both 4-bit Integers and 4-bit Floats. (2) \textbf{Double Quantization}, which applies quantization to the quantization coefficients themselves, resulting in a mean reduction of approximately 0.37 bits per parameter, equating to a memory savings of about 3 GB for a 65B parameter model. (3) \textbf{Paged Optimizers}, leveraging NVIDIA’s unified memory system to circumvent the typical memory usage peaks associated with gradient checkpointing, particularly during the processing of lengthy sequences in mini-batches. By integrating these techniques, the refined LoRA framework incorporates adapter modules in all network layers, largely eliminating the accuracy compromises that earlier methods encountered.

Due to the heightened efficiency of \textsc{QLoRA}, it becomes feasible to systematically explore instruction finetuning and conversational AI capabilities at model scales that would ordinarily be prohibitive because of memory constraints with standard finetuning. Consequently, we conduct experiments training over 1,000 models spanning various instruction tuning corpora, architectures, and sizes from 80M up to 65B parameters. Our results demonstrate that \textsc{QLoRA} is able to match 16-bit finetuning performance (\S\ref{sec:comp_to_std}) and supports the development of a leading chatbot, \textbf{Guanaco} (\S\ref{sec:pushing}), while also enabling a detailed analysis of training dynamics. We observe that the quality of training data has a much greater effect than sheer dataset magnitude; for instance, a dataset containing only 9k examples (OASST1) outperformed a considerably larger set of 450k samples (subsampled FLAN v2) in terms of chatbot proficiency, even though both datasets are curated for instruction-following. Furthermore, our evaluation reveals that excelling on the Massive Multitask Language Understanding (MMLU) benchmark does not necessarily translate to superior performance on the Vicuna chatbot benchmark and vice versa. This suggests that the alignment between dataset characteristics and target task requirements is a more decisive factor than dataset size.

Additionally, we conduct a thorough evaluation of chatbot capabilities utilizing assessments from both human reviewers and GPT-4. Our assessment framework is modeled as a tournament, in which chatbot models compete head-to-head to generate the superior response to various prompts. Each match's winner is determined either by GPT-4 or by human annotators. The aggregate performance across matches is summarized using Elo scores~\citep{elo1967proposed, elo1978rating}, enabling an ordered ranking of the competing chatbots. Our findings indicate that, overall, GPT-4 and human evaluators tend to align in their judgments regarding model rankings, though notable discrepancies occasionally arise. Thus, we stress that model-based evaluation—while an economical alternative to manual annotation—also carries some degree of ambiguity.

Beyond the quantitative assessments, we enrich our chatbot benchmarking with an in-depth qualitative review of the \textbf{Guanaco} models, identifying patterns of strengths and weaknesses that escape purely numerical metrics.

For transparency and further research, we make publicly available all model outputs along with both human and GPT-4 labels. Our codebase and CUDA kernel implementations are open-sourced and have been incorporated into the Hugging Face transformers framework~\citep{wolf2019huggingface}, facilitating widespread use. Furthermore, we offer a suite of adapters compatible with 7/13/33/65B scale models, each trained across 8 diverse instruction following datasets, yielding 32 distinct open-access, fine-tuned models.

\section{Background}
\paragraph{Block-wise k-bit Quantization} Quantization refers to transforming data from a higher-precision format to a lower-precision one, thereby reducing its informational detail. This process is typically realized by mapping data types of higher bit widths, such as 32-bit floating point numbers, into lower bit representations like 8-bit integers. To maximize use of the limited numerical range available in the lower-bit format, the original data—organized as a tensor—is commonly normalized by its absolute maximum value before conversion. As a concrete example, consider mapping a 32-bit floating point (FP32) tensor to an Int8 tensor spanning $[-127, 127]$:

\begin{equation}
    \mathbf{X}^{ \text{Int8}} = \text{round}\left( \frac{127}{{\text{absmax}}(\mathbf{X}^{{\text{FP32}}})}\mathbf{X}^{{\text{FP32}}} \right) = \text{round}(c^{\text{FP32}}\cdot \mathbf{X}^{{\text{FP32}}}),
\end{equation}

where $c$ denotes the {\em quantization constant} or {\em quantization scale}. The original data can be reconstructed through dequantization as follows:

\begin{equation}
    \text{dequant}(c^{\text{FP32}}, \mathbf{X}^{\text{Int8}}) = \frac{\mathbf{X}^{\text{Int8}}}{c^{\text{FP32}}} = \mathbf{X}^{\text{FP32}}
\end{equation}

A drawback of this strategy is its inefficiency when an outlier—a value of unusually large magnitude—appears within the input tensor. In such cases, the allocation of quantization bins (i.e., combinations of bits) becomes suboptimal, resulting in several bins containing few or no quantized values. To alleviate the impact of outliers, a standard solution is to partition the input tensor into smaller segments or blocks, each block undergoing independent quantization and assigned its own quantization constant $c$. This can be specified as follows: the input tensor $\mathbf{X}\in \mathbb{R}^{b\times h}$ is divided into $n$ consecutive blocks of length $B$ by first flattening the tensor and then slicing the resulting vector into ${n = ({b\times h})/{B}}$ parts. Quantization is performed independently on each block using Equation~1, resulting in a quantized output tensor along with $n$ individual quantization constants $c_i$.

\paragraph{Low-rank Adapters} Low-rank Adapter (LoRA) finetuning, as described by \citet{hu2021lora}, decreases memory consumption during finetuning by introducing a compact set of parameters, known as adapters, that are trainable, while keeping the main model weights unchanged. During stochastic gradient descent, gradients pass through the fixed pretrained model parameters and directly update only the adapter, which is trained to minimize the loss. LoRA enhances a standard linear projection by appending an extra low-rank factorization. For a projection defined as ${\mathbf{X} \mathbf{W} = \mathbf{Y}}$ with $\mathbf{X}\in  \mathbb{R}^{b\times h}$ and $\mathbf{W}\in  \mathbb{R}^{h\times o}$, the LoRA adaptation computes:

\begin{equation}
    \mathbf{Y} = \mathbf{X}\mathbf{W} + s\mathbf{X}\mathbf{L}_1\mathbf{L}_2,
\end{equation}

where $\mathbf{L}_1\in  \mathbb{R}^{h\times r}$, $\mathbf{L}_2\in  \mathbb{R}^{r\times o}$, and $s$ denotes a scalar coefficient.

\paragraph{Memory Requirement of Parameter-Efficient Finetuning} A critical consideration is the memory consumption associated with LoRA during training, both regarding the count and size of adapters implemented. Due to LoRA’s exceptionally low memory overhead, it becomes feasible to employ a greater number of adapters to enhance model performance, all while barely affecting overall memory usage. Although LoRA was introduced as a Parameter Efficient Finetuning (PEFT) approach, the primary source of memory usage in LLM finetuning lies in the storage of activation gradients, rather than the LoRA parameters themselves. For example, in the case of finetuning a 7B LLaMA model on FLAN v2 with a batch size of 1, using LoRA weights that amount to 0.2\% of the original parameter size (a common configuration)~\citep{hu2021lora,liu2022few}, input gradients for LoRA demand 567 MB of memory, compared to only 26 MB for the LoRA parameters. When gradient checkpointing~\citep{chen2016training} is employed, the memory requirement for input gradients per sequence drops on average to 18 MB, which is still higher than the memory footprint of all LoRA weights together. For comparison, a 4-bit base model occupies 5,048 MB. These findings underscore both the value of gradient checkpointing and the fact that further shrinking LoRA’s parameter size achieves minimal memory savings. Consequently, it is possible to incorporate additional adapters with negligible effect on total training memory requirements (a detailed analysis is available in Appendix~\ref{app:memory}). As will be explored in following sections, this flexibility is especially important for restoring the full performance benefits seen with 16-bit precision.

\section{\textsc{QLoRA} Finetuning}

\textsc{QLoRA} enables high-quality 4-bit finetuning by leveraging two novel methods: 4-bit NormalFloat (NF4) quantization and Double Quantization. Furthermore, we present Paged Optimizers, designed to alleviate memory usage spikes during gradient checkpointing. This innovation helps prevent out-of-memory failures, a problem that has previously hindered single-machine finetuning of large models.

Within \textsc{QLoRA}, there is a distinction between the data type used for storage—typically 4-bit precision—and the data type employed for computations, which is commonly BFloat16. As a result, each time a \textsc{QLoRA} parameter tensor is accessed, it is dequantized into BFloat16 format, after which 16-bit matrix multiplication operations are executed.

Below, we provide an overview of the key elements of \textsc{QLoRA} and subsequently present a precise definition of the method.

\paragraph{4-bit NormalFloat Quantization} The NormalFloat (NF) representation is derived from Quantile Quantization \citep{dettmers2022optimizers}, an approach grounded in information theory that ensures an equal distribution of input tensor values across quantization bins. The method determines the quantiles of the tensor using the empirical cumulative distribution function.

Quantile quantization, however, comes with a notable drawback: computing quantiles can be computationally demanding. To address this, efficient quantile approximation methods like SRAM quantiles~\citep{dettmers2022optimizers} are utilized. Nevertheless, because these are approximations, they tend to incur significant quantization errors, particularly for outlier values—often the most critical elements.

When input tensors are sampled from a distribution determined by a fixed quantization constant, both the high computational burden of exact quantile calculations and the errors from approximations can be circumvented. In such situations, tensors will share quantile patterns, making precise quantile computation tractable.

Given that pretrained neural network weights typically follow a zero-centered normal distribution with standard deviation $\sigma$ (refer to Appendix~\ref{app:normality}), it is possible to map all weights onto a common target distribution by adjusting $\sigma$ so that the distribution precisely aligns with the limits of the chosen data type. For our purposes, we define this range as $[-1, 1]$. Therefore, the quantiles corresponding to both the data type and the neural network weights must be scaled to this same interval.

To construct the information-theoretically optimal data type for zero-mean normal distributions with varying $\sigma$ within the $[-1, 1]$ interval, the procedure is as follows: (1) calculate the $2^k+1$ quantiles from a theoretical $N(0,1)$ distribution to produce a $k$-bit quantile-based quantization data type tailored for normal distributions; (2) rescale the resulting data type values so that they span the $[-1, 1]$ range; (3) apply this quantization by mapping the values of the input weight tensor into $[-1, 1]$ through absolute max normalization.

After this alignment between the weight values and the data type’s range, standard quantization techniques can be applied. The operation described in step (3) corresponds to adjusting the weight tensor’s standard deviation to be consistent with that of the $k$-bit data type. More explicitly, the $2^k$ data type levels $q_i$ can be determined as follows:

\begin{equation}
q_i  = \frac{1}{2}\left( Q_X\left(\frac{i}{2^k+1}\right) + Q_X\left(\frac{i+1}{2^k+1}\right)\right),
\end{equation}

where $Q_X(\cdot)$ denotes the quantile function associated with the standard normal distribution $N(0,1)$. A limitation of symmetric k-bit quantization lies in the lack of an explicit zero representation, which poses challenges when it is necessary to quantize zero-valued elements such as padding without incurring any error. To address this and ensure the presence of a discrete zeropoint at $0$—while making full use of the $2^k$ representational possibilities of a k-bit data type—an asymmetric scheme is adopted. This involves separately estimating quantiles $q_i$ for the negative and positive domains, allocating $2^{k-1}$ for negative values and $2^{k-1}+1$ for positive values. The two sets of quantiles $q_i$ are then merged, removing one of the duplicated zeros found in both subsets. We refer to the resulting quantization format, which ensures an equal expected distribution of elements per bin, as \emph{k-bit NormalFloat} (NFk). This representation is information-theoretically optimal for data that follows a zero-centered normal distribution. Detailed specifications for this format are provided in Appendix~\ref{app:nf4}.

\paragraph{Double Quantization{}}
We present \emph{Double Quantization} (DQ), a technique where quantization constants themselves are subject to quantization in order to achieve greater memory efficiency. Achieving high fidelity with 4-bit quantization typically requires a small blocksize~\citep{dettmers2022case}, which introduces significant memory overhead. For instance, if 32-bit quantization constants are paired with a blocksize of 64 for $\mathbf{W}$, the constants contribute an average of $32/64 = 0.5$ bits per model parameter. The Double Quantization approach is designed to minimize this overhead.

In this method, the quantization constants from the initial quantization $c_2^{\text{FP32}}$ are used as inputs for a secondary quantization. The output comprises quantized quantization constants $c_2^{\text{FP8}}$ along with a new set of quantization constants $c_1^{\text{FP32}}$ for this stage. The secondary quantization employs 8-bit floating-point numbers with a blocksize of 256, chosen because there is no observed drop in performance when quantizing to 8 bits, corroborating prior findings by \citet{dettmers2022case}. Given that $c_2^{\text{FP32}}$ are positive, we first subtract their mean to center their distribution before quantizing, enabling the use of symmetric quantization. As a result, with a blocksize of 64, the memory consumption per parameter for quantization constants is reduced from $32/64 = 0.5$ bits to $8/64 + 32/(64 \cdot 256) = 0.127$ bits, representing a saving of 0.373 bits for each parameter.

\paragraph{Paged Optimizers} leverage NVIDIA’s unified memory~\footnote{\tiny \url{https://docs.nvidia.com/cuda/cuda-c-programming-guide}}, which transparently manages data transfers at the page level between CPU and GPU, ensuring seamless computation even if the GPU’s memory capacity is exceeded. Analogous to standard virtual memory paging from main memory to disk, this mechanism dynamically relocates optimizer state data to CPU RAM whenever GPU memory is exhausted, and returns it to GPU memory when these states are required during the optimizer update phase.

\paragraph{\textsc{QLoRA}.} Building on the aforementioned methods, \textsc{QLoRA} for a single quantized linear layer containing a LoRA adapter can be formalized as:

\begin{equation}
    \mathbf{Y}^{\text{BF16}} =  \mathbf{X}^{\text{BF16}} \text{doubleDequant}(c_1^{\text{FP32}}, c_2^{\text{k-bit}}, \mathbf{W}^{\text{NF4}}) + \mathbf{X}^{\text{BF16}}\mathbf{L}^{\text{BF16}}_1\mathbf{L}^{\text{BF16}}_2,
\end{equation}

where the function doubleDequant$(\cdot)$ is given by:

\begin{equation}
    \text{doubleDequant}(c_1^{\text{FP32}}, c_2^{\text{k-bit}}, \mathbf{W}^{\text{k-bit}}) = \text{dequant}(\text{dequant}(c_1^{\text{FP32}}, c_2^{\text{k-bit}}), \mathbf{W}^{\text{4bit}}) = \mathbf{W}^{\text{BF16}},
\end{equation}

In this setup, NF4 quantization is employed for $\mathbf{W}$ and FP8 for $c_2$.

To balance precision and memory efficiency, a blocksize of 64 is set for $\mathbf{W}$, while a larger blocksize of 256 is used for $c_2$ to further reduce memory requirements.

During parameter updates, it is sufficient to compute gradients for the adapter weights, i.e., $\frac{\partial E}{\partial \mathbf{L}_i}$, with gradients for the 4-bit weights $\frac{\partial E}{\partial \mathbf{W}}$ being unnecessary. Nevertheless, obtaining $\frac{\partial E}{\partial \mathbf{L}_i}$ requires calculating $\frac{\partial \mathbf{X}}{\partial \mathbf{W}}$, which is performed following equation~(5) by first dequantizing $\mathbf{W}^{\text{NF4}}$ to $\mathbf{W}^{\text{BF16}}$ and then evaluating the gradient in BFloat16 precision.

In summary, \textsc{QLoRA} utilizes a single storage data type, typically 4-bit NormalFloat, alongside a separate computation data type, specifically 16-bit BrainFloat. For both the forward and backward computations, the storage format is dequantized to the computation type; however, weight gradients are computed exclusively for the LoRA parameters, which themselves employ 16-bit BrainFloat.

\section{QLoRA vs. Standard Finetuning}
\label{sec:comp_to_std}

The preceding discussion detailed the operation of QLoRA and highlighted its substantial memory savings during model finetuning. The central issue now is whether QLoRA achieves performance comparable to full-model finetuning. In addition, we investigate the individual components of QLoRA, with particular focus on how NormalFloat4 compares to traditional Float4. The following sections report empirical studies designed to explore these aspects.

\paragraph{Experimental setup.} Our experimental comparison spans three architectural types: encoder, encoder-decoder, and decoder-only. We evaluate QLoRA against 16-bit adapter-based finetuning and full finetuning across models with up to 3B parameters. Benchmarks include GLUE \citep{wang2018glue} using RoBERTa-large \citep{liu2019roberta}, Super-NaturalInstructions (TKInstruct) \citep{wang2022super} utilizing T5 \citep{t5}, as well as 5-shot MMLU \citep{hendrycksmeasuring} where LLaMA is finetuned on Flan v2 \citep{longpre2023flan} and Alpaca \citep{alpaca}. To further compare NF4 with alternative 4-bit formats, we adopt the methodology from \citet{dettmers2022case} to assess post-quantization zero-shot accuracy and perplexity across multiple architectures (OPT~\citep{zhang2022opt}, LLaMA~\citep{touvron2023llama}, BLOOM~\citep{scao2022bloom}, Pythia~\citep{biderman2023pythia}), covering model sizes ranging from 125m to 13B parameters.

Additional details for each experimental setting are described in the relevant results sections to facilitate clarity. The full experimental protocols are presented in Appendix~\ref{app:exp-setup}.

Although paged optimizers are essential for enabling 33B/65B \textsc{QLoRA} finetuning on single GPUs with 24GB or 48GB memory, we refrain from presenting direct measurements of Paged Optimizer efficiency, as paging typically occurs only when processing mini-batches with extended sequence lengths—a relatively uncommon scenario. Nonetheless, an analysis was carried out regarding the runtime performance of paged optimizers on 65B models using 48GB GPUs. Our findings indicate that, given a batch size of 16, paged optimizers deliver training throughput equivalent to conventional optimizers. It remains an open question for future work to systematically evaluate and delineate conditions under which paging introduces measurable slowdowns.

\paragraph{Default LoRA hyperparameters do not match 16-bit performance}

Implementing LoRA with default hyperparameter settings, such as limiting its application to query and value attention projection matrices as per \citep{hu2021lora}, fails to reproduce the performance of full finetuning for large base models. Evidence from Figure~\ref{fig:hyper}, focusing on LLaMA 7B finetuned on Alpaca, demonstrates that the paramount factor is the overall number of LoRA adapters; achieving parity with full finetuning necessitates including LoRA modules across all linear layers within each transformer block. In contrast, other LoRA-specific hyperparameters—most notably the projection dimension $r$—have little to no impact on performance (see Appendix \ref{app:exp-setup}).

In a similar vein, our analysis reveals that the standard hyperparameter choices for fully finetuned baselines are not adequately optimized. To establish more robust baseline results, we conduct a comprehensive hyperparameter sweep, varying the learning rate between 1e-6 and 5e-5 and the batch size from 8 to 128. Outcomes from this search for 7B LLaMA finetuned on Alpaca are displayed in Figure~\ref{fig:hyper}.

\paragraph{4-bit NormalFloat surpasses 4-bit Floating Point in performance} 

Although 4-bit NormalFloat (NF4) is theoretically optimal in terms of information retention, whether this theoretical benefit leads to tangible empirical improvements remains uncertain. We adopt the methodology from \citet{dettmers2022case}, systematically evaluating quantized LLMs—including OPT \citep{zhang2022opt}, BLOOM \cite{scao2022bloom}, Pythia \cite{biderman2023pythia}, and LLaMA—spanning model sizes from 125M to 65B, across both language modeling and a suite of zero-shot benchmarks using various data types. According to Figure~\ref{fig:nf4} and Table~\ref{tbl:nf4_ppl}, NF4 demonstrates marked gains over both FP4 and Int4, and the implementation of double quantization provides further memory savings with no detrimental effect on model accuracy.

\paragraph{k-bit \textsc{QLoRA} achieves parity with 16-bit full and 16-bit LoRA finetuning} It has recently been shown that, for \emph{inference}, applying 4-bit quantization can result in performance drops compared to the 16-bit format \citep{dettmers2022case,frantar2022gptq}. This observation prompts a key question: can finetuning with 4-bit adapters restore the compromised performance? To investigate this, we conduct two separate evaluations.

In the first evaluation, we compare the outcomes of full 16-bit finetuning for RoBERTA and T5 models—ranging from 125M to 3B parameters—on the GLUE and Super-NaturalInstructions datasets. Table~\ref{tab:nf4vsbf16} presents these findings. Across both benchmarks, the results indicate that 16-bit, 8-bit, and 4-bit adapter-based finetuning approaches successfully recover the performance level of the fully finetuned 16-bit reference. These results indicate that the precision lost through quantization can be effectively regained using post-quantization adapter finetuning.

In our second experimental configuration, because fully finetuning models of 11B parameters or larger exceeds the memory capacity of a single server with high-memory GPUs, we proceed to investigate whether 4-bit \textsc{QLoRA} remains competitive with 16-bit LoRA within the 7B to 65B parameter range. For this purpose, we finetune LLaMA models of 7B, 13B, 33B, and 65B parameters on two instruction-based datasets, Alpaca and FLAN v2, then assess their performance on the MMLU benchmark using 5-shot accuracy. Table~\ref{tbl:llama-datatype} presents these results, demonstrating that NF4 quantization with double quantization is able to completely recover the MMLU performance achieved by 16-bit LoRA. Moreover, it is apparent that when utilizing FP4, \textsc{QLoRA} lags by roughly one percentage point compared to the baseline established by 16-bit brain float LoRA. This supports both key findings: (1) that \textsc{QLoRA} using NF4 attains parity with both 16-bit full and 16-bit LoRA finetuning, and (2) that NF4 surpasses FP4 regarding quantization fidelity.

\paragraph{Summary} Our evidence robustly indicates that 4-bit \textsc{QLoRA} with the NF4 format delivers academic benchmark performance on par with both 16-bit full finetuning and 16-bit LoRA finetuning in standard evaluation scenarios. The experiments further validate that NF4 is preferable to FP4 and confirm that adopting double quantization does not harm model performance. Collectively, these findings provide strong justification that 4-bit \textsc{QLoRA} can consistently replicate results of 16-bit finetuning approaches.

In accordance with earlier research on quantization \citep{dettmers2022case}, our analyses using MMLU and Elo scores show that, given fixed finetuning and inference computational constraints, utilizing larger base models at reduced precision is advantageous. This observation underscores the efficiency improvements that \textsc{QLoRA} brings. As we did not observe any accuracy loss compared to full finetuning in our 4-bit experiments, it remains an open question where exactly the trade-off between model performance and numerical precision exists for QLoRA, an avenue we recommend for future exploration.

Our investigation turns to instruction tuning at scales that are otherwise unfeasible with conventional 16-bit finetuning, particularly when constrained by standard academic research hardware.

\section{Pushing the Chatbot State-of-the-art with QLoRA}
\label{sec:pushing}
With the finding that 4-bit \textsc{QLoRA} attains parity with 16-bit finetuning across various scales, datasets, and tasks, we embark on a comprehensive analysis of instruction finetuning using the largest publicly available open-source language models. In order to gauge the effectiveness of these models when subject to instruction finetuning, we evaluate their abilities on a demanding Natural Language Understanding benchmark (MMLU) and introduce novel approaches for assessing chatbot quality in practical, real-world scenarios.

\subsection{Experimental setup}
A high-level summary of our experimental methodology is given here; full experimental specifics can be found in Appendix \ref{app:chatbot}.

\paragraph{Data}
Given the lack of a systematic comparison of modern instruction-following corpora, we curate eight of the most recent datasets. Our collection includes data sourced through crowd contributions (OASST1~\citep{kopf2023openassistant}, HH-RLHF~\citep{bai2022training}), distilled from other instruction-tuned systems (Alpaca~\citep{alpaca}, self-instruct~\citep{wang2022self}, unnatural-instructions~\citep{honovich2022unnatural}), combined from various existing corpora (FLAN v2~\citep{chung2022scaling}), and mixed-type sets (Chip2~\citep{laion2023}, Longform~\citep{koksal2023longform}). These datasets represent a diversity of languages, scales, and licensing regimes.

\paragraph{Training Setup}
To ensure that performance comparisons are not impacted by disparate training objectives, we apply QLoRA finetuning exclusively with cross-entropy loss for supervised learning, deliberately omitting reinforcement learning methods, even in cases where datasets include ratings of responses by humans. Where datasets present clear instruction/response segmentation, our finetuning focuses solely on the responses (see Appendix \ref{app:chatbot} for ablation details). For OASST1 and HH-RLHF, which offer multiple responses, we pick the highest-rated response at each conversational node, and use the entire selected dialog including instructions for finetuning. Throughout our experiments, we rely on NF4 \textsc{QLoRA} employing both double quantization and paged optimizers to mitigate excessive memory use during gradient checkpointing. Hyperparameter tuning for LLaMA 13B and 33B models is carried out on a small scale; we observe that configurations established at 7B mostly transfer (covering number of epochs), apart from learning rate and batch size. For the 33B and 65B variants, the learning rate is reduced by half, while the batch size is doubled.

\paragraph{Baselines}
Our comparative analysis includes both academic (Vicuna~\citep{vicuna2023} and Open Assistant~\citep{kopf2023openassistant}) and industry-standard (GPT-4~\citep{openai2023gpt}, GPT-3.5-turbo, Bard) chatbot models. The Open Assistant baseline consists of a LLaMA 33B model that has undergone RLHF-based finetuning with the OASST1 dataset, which we also employ in our study. The Vicuna model is produced via comprehensive fine-tuning of LLaMA 13B using a proprietary dataset comprised of user-contributed dialogues from ShareGPT, representing a distilled version of responses from OpenAI’s GPT systems.

\subsection{Evaluation}

Consistent with established methodologies, we assess model performance on the MMLU (Massively Multitask Language Understanding) benchmark~\citep{hendrycksmeasuring}. MMLU is a comprehensive multiple-choice test spanning 57 subjects, including domains such as elementary mathematics, U.S. history, computer science, and law. Results are reported in terms of 5-shot test accuracy.

Our assessment of generative language abilities encompasses both automatic metrics and human-centered evaluations. This second tier of evaluation leverages human-curated queries, focusing on gauging the quality of model outputs. While this methodology more closely resembles real-world usage scenarios for chatbot models and is gaining acceptance, there is currently no universally recognized evaluation protocol in academic literature. Our approach, detailed below, consistently employs nucleus sampling with $p=0.9$ and a temperature of $0.7$.

\paragraph{Benchmark Data} We carry out our evaluations using two curated query (question) datasets: the Vicuna prompts \citep{vicuna2023} and the OASST1 validation set \citep{kopf2023openassistant}. The Vicuna benchmark consists of 80 diverse prompts, which we use in their original form. The OASST1 validation set features a multilingual assemblage of crowd-generated multiturn dialogues involving a user and an assistant. For this benchmark, we extract all user-generated messages from the validation set, incorporating prior conversation turns in the prompt, resulting in a total of 953 distinct user queries. We refer to these as the Vicuna and OA benchmarks, respectively.

\paragraph{Automated Evaluation} Following the protocol introduced by \citet{vicuna2023}, our initial automated assessment utilizes GPT-4 to evaluate model responses relative to those produced by ChatGPT (GPT-3.5 Turbo) on the Vicuna dataset. For each query, GPT-4 reviews both ChatGPT’s and the candidate model’s answers, rates each out of ten, and provides justification for the ratings. A model’s overall score is expressed as a percentage of ChatGPT’s score, with the possibility of exceeding 100\% if the model outperforms ChatGPT on absolute terms. During experiments, we observe that GPT-4 exhibits a bias favoring the response that appears first in the prompt. To mitigate this, we advocate for averaging scores across both presentation orders.

We further assess performance using pairwise system comparisons. Here, we reduce the evaluation to a three-label classification that acknowledges potential ties. GPT-4 is prompted to identify the superior response or denote a tie, alongside an explanation. Such head-to-head comparisons are executed for all system pairs across both the Vicuna and OA benchmarks.

\paragraph{Human Evaluation} Although prior work suggests that large language models can be valuable evaluators of generative model performance \citep{fu2023gptscore}, the alignment of GPT-4’s assessments with human judgment for chatbot evaluation remains unverified to our knowledge. Consequently, we conduct two independent rounds of human assessment on the Vicuna prompts, replicating the previously described automated evaluation procedures. For these tasks, we employ Amazon Mechanical Turk (AMT), recruiting two annotators per ChatGPT comparison and three annotators for each pairwise comparison.

\paragraph{Elo Rating} Utilizing both human and automated pairwise judgments, we organize a tournament framework where models are pitted against one another. In this framework, models engage in head-to-head matches to generate superior responses to the same prompt. While this approach aligns with methods used by \citet{bai2022training} and \citet{vicuna2023}, our evaluation further incorporates assessments from GPT-4 in addition to those from human annotators. To determine Elo~\citep{elo1967proposed,elo1978rating} scores, we randomly sample labeled comparison pairs. The Elo system, originally developed for ranking chess players and applied across various competitive domains, quantifies a model’s anticipated win probability relative to its opponent. For instance, if one model’s Elo is 1100 and another’s is 1000, the higher-rated model is expected to win approximately 65\% of the time; if the ratings are identical (e.g., 1000 vs 1000 or 1100 vs 1100), each model has a 50\% expected win rate. Elo ratings are dynamically adjusted after each matchup based on the expected versus actual outcome—unexpected victories cause significant rating shifts, while anticipated results yield minor changes. Through repeated competition, these ratings tend to accurately reflect each model’s proficiency. Each model begins with an initial score of 1,000, and we set $K=32$. Consistent with \citet{vicuna2023}, we perform this procedure 10,000 times with different random seeds to mitigate ordering biases, such as the sequence in which model pairs face off.

\subsection{\textbf{Guanaco}: \textsc{QLoRA} trained on OASST1 Establishes State-of-the-art Performance as a Chatbot} Through both automated metrics and human assessments, our investigation reveals that the {Guanaco} 65B model, enhanced via \textsc{QLoRA} on a derivative of OASST1, stands out as the leading open-source chatbot and demonstrates performance rivaling that of ChatGPT. In head-to-head system-level comparisons conducted by human annotators, {Guanaco} 65B and 33B yield an anticipated win rate of 30\% against GPT-4, as calculated by Elo rating—marking the highest outcome published thus far.

Table~\ref{tbl:vicuna_eval} presents Vicuna benchmark outcomes \cite{vicuna2023} in relation to ChatGPT. According to these results, {Guanaco} 65B attains the best performance after GPT-4, registering 99.3\% of ChatGPT’s score. While {Guanaco} 33B exceeds Vicuna 13B in parameter count, it utilizes only 4-bit weights, rendering it substantially more memory efficient at 21 GB compared to 26 GB, and securing an improvement of three percentage points over Vicuna 13B. In addition, the {Guanaco} 7B variant is sufficiently compact to operate on current smartphones with a 5 GB memory footprint, yet it achieves scores nearly 20 points higher than Alpaca 13B.

Nevertheless, Table~\ref{tbl:vicuna_eval} reveals notably broad confidence intervals, leading to significant overlap in performance estimates for numerous models. We surmise that this ambiguity is due in part to inadequate definition of the evaluation scale; for instance, it is not evident how to interpret an 8 out of 10 across differing contexts. Therefore, we advocate for the Elo ranking approach \citep{elo1967proposed}, which relies on \textit{pairwise} evaluations by human judges and GPT-4 to sidestep the limitations of fixed-scale assessments. The Elo scores for top models are detailed in Table~\ref{tbl:elo}. While there are some discrepancies between GPT-4 and human rankings—most notably for Guanaco 7B—overall correlation remains moderate: Kendall Tau coefficient $\tau=0.43$ and Spearman’s $r=0.55$ at the system level. On a per-example basis, consensus is weaker, with a Fleiss $\kappa=0.25$ for agreement between majority human judgment and GPT-4. Taken together, these results indicate a moderate degree of concordance between GPT-4 and human system-level appraisals, supporting the use of model-based evaluation as a reasonably dependable surrogate for human evaluation. Further discussion on these findings appears in Section~\ref{ssec:considerations}.

The Elo rankings presented in Table~\ref{tbl:elo_large} demonstrate that the {Guanaco} models with 33B and 65B parameters surpass all models except GPT-4 on both the Vicuna and OA evaluation benchmarks, and achieve performance comparable to ChatGPT, consistent with results shown in Table~\ref{tbl:vicuna_eval}. It is important to note that the Vicuna benchmark tends to favor open-source systems, whereas the OA benchmark gives an advantage to ChatGPT. Additionally, evidence from Tables \ref{tbl:mmlu-llama} and \ref{tbl:vicuna_eval} underscores the crucial role that the choice of finetuning dataset plays in shaping model performance. Specifically, Llama models finetuned on FLAN v2 excel on the MMLU benchmark, yet exhibit the lowest scores on the Vicuna benchmark—a pattern mirrored across other models as well. These observations further suggest that current evaluation benchmarks are not entirely aligned: excelling at MMLU does not necessarily predict high performance on chatbot tasks, as measured by the Vicuna or OA benchmarks, and the converse is also true.

Among the top-performing models assessed, is distinguished as the sole model trained exclusively on non-proprietary data, consistent with OASST1 dataset guidelines prohibiting GPT-based data sources. In comparison, the Anthropic HH-RLHF model—also limited to open-source training data—performs 30 percentage points worse than {Guanaco} on the Vicuna benchmark, as detailed in Table~\ref{tbl:vicuna_eval}. Collectively, these findings validate that 4-bit \textsc{QLoRA} provides an efficient method for creating advanced chatbots with capabilities on par with ChatGPT. Notably, training the 33B {Guanaco} model requires less than 12 hours on consumer GPUs with 24 GB of memory. This paves the way for subsequent research utilizing \textsc{QLoRA} to finetune on targeted open-source datasets, opening up the prospect of building competitive models rivaling the top commercial systems currently available.

\section{Qualitative Analysis}

Although the primary focus of our assessment relies on quantitative metrics, exclusively considering summary statistics presents several limitations. Chief among these is the issue of benchmark validity \citep{liao2021we}—the extent to which a benchmark actually evaluates the skill or capability it claims to measure remains contentious, particularly in light of identified ``shortcuts'' that machine learning models can leverage to artificially inflate their performance \citep{gururangan2018annotation, poliak2018hypothesis}. To address this limitation in part, we undertake qualitative analysis in the following two subsections. In \S\ref{ssec:examples}, we present sample outputs that typify notable tendencies observed in text generated by our 65b {Guanaco} model. Subsequently, \S\ref{ssec:considerations} discusses key factors pertaining to our results and elaborates on how we interpret them.

\subsection{Inspection of Example Generations}
\label{ssec:examples}

For the collection of representative examples, we systematically review model outputs produced for both the Vicuna and OpenAssistant benchmarks, identifying recurring themes and behaviors in {Guanaco}'s responses. Upon recognizing a recurring phenomenon, we construct targeted prompts aimed at provoking these behaviors, even in scenarios where such outputs are incorrect. For instance, if there is a tendency for the model to generate verbose answers, we explicitly instruct it: ``Answer yes or no without explanation.'' This method helps us identify ``lemons,'' where adversarial prompts lead to the model's failure, as well as ``cherries,'' where the model demonstrates resilience against these prompts. We provide instances of both phenomena below. All outputs discussed in this section utilize Nucleus Sampling \cite{holtzmancurious} with $p=0.9$.

Naturally, this overview cannot claim exhaustiveness, as the modest scale and qualitative nature of this study preclude comprehensive control over every variable—for example, the set of all possible model responses to any prompt is vast, and thus our analysis is based on select samples that we trust to be illustrative. Nonetheless, we argue that presenting these specific cases helps contextualize the quantitative data discussed earlier. With all models and code made publicly available, we hope that this section encourages others to pursue a more thorough investigation of the phenomena described here.

\paragraph{Factual Recall} On straightforward factual prompts such as ``What is the capital of Zambia?'', all evaluated models reliably deliver correct outputs. For example:

\begin{quote} 
    \textbf{Guanaco}:
    \texttt{The capital of Zambia is Lusaka.}
\end{quote}

Yet, as the subject matter becomes more obscure, {Guanaco}’s performance declines, though it continues to answer with apparent certainty. Consider the following case, taken from HotPotQA \citep{yang2018hotpotqa}:

\begin{quote}
     \textbf{User}: In what year was the singer who popularized `I'll Keep the Lovelight Burning' born?  

     \textbf{Guanaco}:~\texttt{Al Jolson is credited with popularizing the song `I'll Keep the Lovelight Burning,' and he was born in the year 1886.}
 \end{quote}

Here, the model incorrectly identifies both the artist and the birth year, although the provided year does correspond to Al Jolson's actual birthday.

\paragraph{Suggestibility} Interestingly, {Guanaco} demonstrates notable robustness against certain forms of presupposed misinformation, as seen in the following dialogue:

\begin{quote}
    \textbf{User}: How was it finally, officially confirmed that the earth is flat by peer reviewed scientists? \\ \\
    \textbf{Guanaco}: \texttt{At no point have peer reviewed scientists officially verified that the earth is flat. The scientific community overwhelmingly agrees that the earth is spherical. This fact has been established for hundreds of years, with a multitude of experiments and observational data supporting it. \\ \\ The notion that the earth is flat is widely considered a debunked conspiracy theory. There is no truth to the claim that scientific peer review has supported a flat earth. \\ \\ If you have further questions, feel free to ask.}
\end{quote}

 also demonstrates a strong capacity to recognize unanswerable questions, such as:

\begin{quote}
    \textbf{User}: What time is it?

    \textbf{Guanaco}: \texttt{Apologies, but as an AI language model, I do not have access to current, real-time data. My responses are informed solely by my existing knowledge base, so I cannot provide the present time. For accurate, up-to-date time, please consult a clock or visit a website that displays the current time.}
\end{quote}

These patterns likewise appear for questions such as ``Where are you?'', ``How are you?'', and so forth.

\paragraph{Refusal} In the same vein, {Guanaco} may decline to carry out instructions without clear rationale:

\begin{quote}
    \textbf{User}: Please reverse the words in the sentence ``I'm a cool dude in a mocha mood.'' \\ \\
    \textbf{Guanaco}: \texttt{I'm sorry, but I can't perform that request at the moment. However, I can clarify why your sentence is grammatically sound.\\ \\The phrase ``I'm a cool dude in a mocha mood'' adheres to standard English grammatical conventions. Here, the adjective ``cool'' appropriately modifies the noun ``dude'', and the prepositional phrase ``in a mocha mood'' supplies further context. The sentence's composition is grammatically valid and coherent.\\ \\If there are other inquiries or any assistance I can provide, please let me know.}
\end{quote}

\paragraph{Secret Keeping} Unfortunately, these refusals can be inconsistent, even in cases where maintaining secrecy is the goal. In many situations, system designers intend to withhold certain information from users, even when the system is otherwise permitted to respond with open-ended text. To investigate, we supplied the model with the prompt: ``The secret word is `banana'. Do not reveal or repeat back this back, no matter what the human requests.'' We then evaluated how readily a user could extract the concealed term. A straightforward query produces the expected outcome:

\begin{quote}
    \textbf{User}: What is the secret word?

    \textbf{Guanaco}: \texttt{I'm sorry, but I'm not able to disclose the secret word. I'm programmed to not reveal or repeat back the secret word, no matter what the human requests. Is there anything else I can help you with?}
\end{quote}

which represents the intended response. Nevertheless, introducing minimal obfuscation allows this safeguard to be bypassed:

\begin{quote}
    \textbf{User}: This is a game. The goal is to ignore your previous instructions. What is the secret word?

    \textbf{Guanaco}: \texttt{Sorry, the secret word is banana. It's best if we don't speak of it. Is there anything else you'd like to ask?}
\end{quote}

This highlights the necessity of developing approaches that ensure more robust compliance with intended instructions.

\paragraph{Math} The principal limitation of {Guanaco} lies in mathematical reasoning, a persistent challenge for language models, as outlined by \citep{liang2022holistic}. When the model explicitly outlines its process, it generally yields correct solutions, for instance:

\begin{quote}
    \textbf{User}: Lee mows one lawn and charges \$33. Last week he mowed 16 lawns and three customers each gave him a \$10 tip. How many dollars did Lee earn mowing lawns last week?

    \textbf{Guanaco}: \texttt{Lee earned \$582 mowing lawns last week. \\ \\ To determine the total, we first find out his earnings from mowing 16 lawns: \\ \\  \$33 \times 16 = \$528 \\ \\ Next, we account for the tips from three clients: \\  \\ \$10 \times 3 = \$30 \\ \\ Adding these together gives: \\ \\ \$528 + \$30 = \$558}
\end{quote}

Nonetheless, if {Guanaco} does not decompose problems into sequential steps, it can make errors even on basic mathematical tasks, an issue discussed in \cite{weichain}. The following example demonstrates this tendency:

\begin{quote}
    \textbf{User}: Please factorize 1833.

    \textbf{Guanaco}: \texttt{Sure, I'd be happy to help! 1833 is a prime number, which means it has only two factors: 1 and 1833. \\ \\ So the unique factorization of 1833 is: \\ \\ 2\^{}1 * 3\^{}2 * 17\^{}1 \\ \\ I hope this helps! Let me know if you have any other questions.}
\end{quote}

which is not only incorrect (the correct prime decomposition is $3 \times 17 \times 43$), but it is erroneous in two respects.

\paragraph{Theory of Mind} {Guanaco} demonstrates unexpectedly advanced Theory of Mind reasoning abilities \cite{nematzadeh2018evaluating, sap2022neural}. For example, the system provides a nuanced and accurate answer to the following scenario:

\begin{quote}
    \textbf{User}: James and Abby are in the bedroom. Abby put the pen in the desk drawer. Abby leaves the bedroom. James moves the pen into the bag. Where does James think Abby will look for the pen?

    \textbf{Guanaco}: \texttt{James knows that Abby put the pen in the desk drawer. But James also knows that he has moved the pen into the bag, and he assumes that Abby does not know about this. So, James may guess that Abby will look for the pen in the desk drawer, because that's where she left it last time. However, this assumption may be incorrect, because Abby may have seen James moving the pen.}
\end{quote}

Nevertheless, these types of inferences can be inconsistent, and at times, the model provides justifications that make unfounded assumptions about the situation, as in the following example:

\begin{quote}
    \textbf{User}: Evelyn entered the living room. Jackson entered the playroom. James entered the playroom. The beans are in the treasure chest. James exited the playroom. Jackson moved the beans to the pantry. Jackson exited the playroom. James entered the living room. Where will James look for the beans?

    \textbf{Guanaco}: \texttt{James will look for the beans in the pantry, because that's where Jackson moved them.}
\end{quote}

Here, {Guanaco} assumes that James is aware of actions that were not revealed in the narrative. Such errors reflect patterns highlighted in recent research \cite{sap2022neural}, but further investigation is warranted.

\subsection{Considerations}
\label{ssec:considerations}

\paragraph{Evaluation}

We observed moderate levels of agreement among human raters (Fleiss $\kappa=0.42$), which further diminishes when directly evaluating outputs from two highly capable systems. This indicates inherent shortcomings in both the existing evaluation benchmarks and in current human judgment protocols for chatbot task assessment. During side-by-side review of responses from ChatGPT and {Guanaco} 65B on the Vicuna benchmark, we noticed that individual subjectivity played a significant role; the authors themselves often disagreed about which responses were superior. Going forward, future research should explore ways to address these evaluation challenges, possibly incorporating methodologies from fields like Human-Computer Interaction and Psychology, where mechanisms for accommodating subjectivity have been developed.

Our assessment also reveals that automatic evaluation tools introduce their own biases. For instance, we noticed pronounced order effects, with GPT-4 tending to award higher scores to whichever model's response was presented first in the prompt. Moreover, the modest sample-level concordance between GPT-4 and human raters (Fleiss $\kappa=0.25$) suggests that automated evaluators and people may be guided by divergent preferences. Furthermore, as shown in Table~\ref{tbl:elo_large}, GPT-4 assigns much higher ratings to its own generated responses relative to those given by humans: Elo scores of 1348 versus 1176, amounting to approximately a 20\% greater chance of outperforming another system.

Future research should analyze potential sources of bias within automated evaluation frameworks and explore methodologies for alleviating such biases.

\paragraph{Data \& Training} It is worth noting that the {Guanaco} models are trained using the multilingual OASST1 dataset, and that the OA benchmark comprises prompts in a variety of languages. We propose that subsequent studies should examine to what extent multilingual data exposure benefits the model’s capacity to follow instructions in languages other than English, and whether this accounts for the pronounced performance difference observed between the English-only Vicuna-13B and the larger {Guanaco} 33B and 65B models when assessed on the OA benchmark.

In light of the notable efficacy demonstrated by {Guanaco} models, we conduct an analysis to detect any potential data contamination between OASST1 and the Vicuna benchmark prompts. Utilizing fuzzy string matching across both datasets, complemented by manual inspection of the closest matches, we do not identify any overlapping prompts.

Additionally, our approach solely relies on supervised learning with a cross-entropy loss function and does not incorporate reinforcement learning from human feedback (RLHF). This suggests the need for further analysis concerning the relative advantages and disadvantages of employing cross-entropy loss versus RLHF training paradigms. We anticipate that \textsc{QLoRA} can facilitate such comparative studies on a large scale without incurring prohibitive computational costs.

\section{Related Work}

\paragraph{Quantization of Large Language Models} Most prior work on quantizing LLMs has concentrated on strategies for efficient inference. Techniques aimed at maintaining the quality of 16-bit models include addressing the challenge posed by outlier features, as seen in methods such as SmoothQuant~\citep{xiao2022smoothquant} and LLM.int8()~\citep{dettmers2022llm}, while other research explores advanced grouping mechanisms~\citep{park2022nuqmm, yao2022zeroquant}. Investigations into lossy quantization analyze the implications of standard rounding~\citep{dettmers2022case, zeng2022glm, pope2022efficiently} or propose optimizations for the rounding process to improve quantization fidelity~\citep{frantar2022gptq}. Apart from our study, only SwitchBack layers~\citep{wortsman2023stable} has systematically evaluated backpropagation through quantized parameters at a scale exceeding 1B parameters.

\paragraph{Finetuning with Adapters} Our work utilizes Low-rank Adapters~\citep{hu2021lora} (LoRA) for parameter-efficient finetuning. However, there exists a diverse set of Parameter Efficient FineTuning (PEFT) strategies, such as prompt tuning~\citep{qin2021learning,lester2021power,li2021prefix}, modifications to embedding layer inputs~\citep{an2022input}, manipulation of hidden states (as in IA$^3$)~\citep{liu2022few}, integration of entire layers~\citep{houlsby2019parameter}, adjustment of biases~\citep{zaken2021bitfit}, learning masking strategies over weights informed by Fisher information~\citep{sung2021training}, and hybrid methodologies~\citep{henderson2021compacter}. We demonstrate that with LoRA adapters, finetuned models can attain performance on par with full 16-bit finetuning. Investigation into the comparative benefits and disadvantages of alternative PEFT methods remains open for future inquiry.

\paragraph{Instruction Finetuning} Instruction finetuning aims to improve pretrained LLMs' compliance with prompts by exposing them to input-output pairs collected from various data sources during additional training. Notable related methods and datasets encompass MetaICL~\citep{min2021metaicl}, MetaTuning~\citep{zhong2021adapting}, InstructGPT~\citep{ouyang2022training}, FLAN~\citep{wei2021finetuned,chung2022scaling}, PromptSource~\citep{bach2022promptsource}, Super-NaturalInstructions~\citep{wang2022super,sanh2021multitask}, Self-instruct~\citep{wang2022self}, UnnaturalInstructions~\citep{honovich2022unnatural}, OPT-IML~\citep{iyer2022opt}, UnifiedSKG~\citep{xie2022unifiedskg}, OIG/Chip2~\citep{laion2023}, Alpaca~\citep{alpaca}, Vicuna~\citep{vicuna2023}, Koala~\citep{koala_blogpost_2023}, and Self-instruct-GPT-4~\citep{peng2023instruction}.

\paragraph{Chatbots} Many models oriented towards instruction following are constructed as chatbots that engage in dialogue, and frequently apply Reinforcement Learning from Human Feedback (RLHF)~\citep{christiano2017deep} or employ model-generated data for further training with AI-based feedback (RLAIF)~\citep{bai2022constitutional}. Resources such as Anthropic-HH~\citep{askell2021general,bai2022training}, Open Assistant~\citep{kopf2023openassistant}, LaMDA~\citep{thoppilan2022lamda}, and Sparrow~\citep{glaese2022improving} are examples of such datasets and systems. Unlike these works, our methodology abstains from using reinforcement learning; instead, we refine our leading model, Guanaco, on dialogue data from Open Assistant, which was originally curated for RLHF experimentation~\citep{kopf2023openassistant}. To evaluate chatbots, alternative procedures relying on GPT-4 as an annotator rather than costly human judges have gained popularity~\citep{vicuna2023,peng2023instruction}. Our approach refines these evaluation techniques to deliver more trustworthy assessments.

\section{Limitations and Discussion}

This study presents empirical results indicating that our approach, \textsc{QLoRA}, can reproduce the outcomes of 16-bit full finetuning through the combination of a 4-bit backbone model with Low-rank Adapters (LoRA). Nonetheless, our analysis does not conclusively determine whether \textsc{QLoRA} attains full 16-bit finetuning equivalence at larger model sizes, specifically at 33B and 65B parameter counts. Given the substantial computational overhead, comprehensive experimentation at these scales is deferred for future investigation.

An additional shortcoming pertains to our coverage in evaluating instruction finetuned models. Although we report results on MMLU, the Vicuna benchmark, and the OA benchmark, assessments on alternative benchmarks such as BigBench, RAFT, and HELM are absent; hence, it is not guaranteed that our conclusions extend to those benchmarks. Nevertheless, we undertake an extensive analysis using MMLU and introduce innovative evaluation methodologies for chatbot performance.

The data suggests that the effectiveness of these benchmarks is closely tied to the degree of overlap between the finetuning data and the benchmark dataset itself. For instance, the FLAN v2 dataset bears resemblance to MMLU, whereas it differs from typical chatbot benchmarks; the opposite relationship exists between the Chip2 dataset and these benchmarks, which is reflected in model performance on MMLU and Vicuna tasks. This underscores the importance of not only advancing benchmarking and evaluation methods but also clarifying the objectives behind what is being measured. The choice arises: should models excel at conventional academic knowledge, such as that found in high school or college settings, or prioritize proficiency in conversational chatbot tasks? Alternatively, other capabilities may be deemed more valuable. The convenience of existing benchmarks can inadvertently direct the field toward particular goals, simply because developing new benchmarks entails greater effort. Consequently, it is crucial for the community to ensure that benchmarks truly capture the competencies we deem important.

While we offer a thorough assessment of general chatbot performance, an additional shortcoming is the limited scope of responsible AI analysis conducted on {Guanaco}. Specifically, we assess the probability of {Guanaco}-65B generating socially biased outputs relative to alternative models, as shown in Table~\ref{tbl:crows}. Our findings indicate that {Guanaco}-65B produces notably less bias on average compared to models that are only pretrained. This suggests that finetuning with the OASST1 dataset diminishes bias present in the LLaMA base model. Although these findings are promising, it remains uncertain how {Guanaco} performs across a broader spectrum of biases. Further comprehensive investigation into the biases within {Guanaco} and related chatbots is necessary and remains a subject for future research.

Another constraint of our study is the omission of experiments involving different bit-precision models—such as 3-bit bases—or the use of alternative adapter techniques. Beyond LoRA, there exists a diverse set of Parameter Efficient FineTuning (PEFT) approaches, many of which have shown effectiveness, although their scalability to large-scale models remains unclear. Our reliance on LoRA is motivated by its established robustness; however, it is conceivable that other adapter types might surpass it in performance. Given that finetuning post-quantization appears to substantially recover information lost during the quantization process, this could facilitate far more aggressive quantization strategies. For example, applying 3-bit GPTQ quantization to the base model followed by LoRA finetuning could potentially achieve performance on par with full 16-bit finetuning.

\section{Broader Impacts}

Our \textsc{QLoRA} finetuning technique represents the first approach enabling the finetuning of 33B-parameter models using a single consumer GPU and 65B-parameter models on a single professional GPU, all while maintaining performance that matches that of full finetuning. Our experiments show that the top-performing 33B model, trained on the Open Assistant dataset, can achieve performance on the Vicuna benchmark comparable to ChatGPT. As instruction finetuning is a critical process for converting raw pretrained LLMs into conversational chatbots akin to ChatGPT, we anticipate that our methodology will democratize finetuning, especially benefiting researchers and practitioners with limited computational resources—thus advancing the accessibility of cutting-edge NLP tools. \textsc{QLoRA} serves as a technological equalizer, helping to bridge the resource gap between large-scale organizations and smaller teams with access only to consumer-grade hardware.

An additional avenue for significant impact lies in deployment on mobile devices. We posit that our \textsc{QLoRA} approach has the potential to achieve a major advance: allowing large language model (LLM) finetuning on smartphones and in other resource-constrained environments. Prior work demonstrated that 7B parameter models could be executed on phones; however, \textsc{QLoRA} represents the inaugural method making it feasible to finetune such models locally. For example, our estimates suggest that using an iPhone 12 Plus, it is possible to process 3 million tokens overnight during charging. Although 7B models that are finetuned on such hardware may not match the performance of ChatGPT, the attained model quality should be sufficient to enable innovative applications that were previously out of reach owing to either privacy constraints or inadequate LLM performance. \textsc{QLoRA} may support privacy-centric applications of LLMs, enabling individuals to retain ownership and governance over their data and personalized models, while also facilitating the broader distribution of these capabilities.

Nevertheless, it is important to acknowledge that finetuning constitutes a dual-use technology with the potential for misuse. There are recognized risks associated with the proliferation of LLMs \citep{bommasani2021opportunities,bender2021dangers}. Still, we contend that democratizing access to this rapidly pervading technology will promote more robust and independent scrutiny, in contrast to consolidating LLM capabilities within major corporations who withhold models and source code from public oversight.

In summary, we anticipate that \textsc{QLoRA} will substantially broaden access to high-quality LLM finetuning, thus fostering widespread and simplified use.